\section{高级并行策略}

\subsection{上下文并行（Context Parallelism）}

上下文并行在序列维度上切分，将长序列分配给不同 GPU 处理。

\begin{figure}[H]
\centering
\includegraphics[width=0.95\textwidth]{slides-images/slide-032.jpg}
\caption{上下文并行：对于 Transformer 的大部分组件（LayerNorm、FFN、残差连接），序列维度是天然独立的\protect\footnotemark}
\end{figure}
\footnotetext{Slides 第32页。}

\begin{itemize}
    \item LayerNorm、FFN、残差连接在序列维度上天然独立，可以直接切分
    \item \textbf{Attention 是难点}：因为需要计算所有 token 对之间的交互
    \item 两种解决方案：Ring Attention（按 block 轮转计算）和 Ulysses Attention（按 head 并行）
\end{itemize}

Llama 3 在长上下文训练阶段（序列长度 131,072）使用了 16-way 上下文并行。

\subsection{流水线并行（Pipeline Parallelism）}

流水线并行将网络的不同层分配给不同 GPU。主要挑战是层间的\textbf{顺序依赖}：

\begin{figure}[H]
\centering
\includegraphics[width=0.95\textwidth]{slides-images/slide-035.jpg}
\caption{朴素流水线并行的``bubble''问题：大部分时间 GPU 在空闲等待\protect\footnotemark}
\end{figure}
\footnotetext{Slides 第35页。}

朴素实现中，$n$-way 流水线并行的利用率仅为 $1/n$（8-way 时最大 MFU 仅 12.5\%）。

\begin{figure}[H]
\centering
\includegraphics[width=0.95\textwidth]{slides-images/slide-037.jpg}
\caption{通过多个 micro-batch 缩小 bubble：4-way 流水线 + 4 个 micro-batch 可达 57\% 理论 MFU\protect\footnotemark}
\end{figure}
\footnotetext{Slides 第37页。}

\begin{importantbox}{缩小 Bubble 的关键：微批次}
通过同时处理多个 micro-batch，让不同 GPU 在同一时刻处理不同 micro-batch 的不同层，可以显著提高 GPU 利用率。micro-batch 越多，bubble 越小，但每个 micro-batch 的激活值都需要存储——又需要激活检查点。
\end{importantbox}

\subsection{张量并行（Tensor Parallelism）}

张量并行将单个权重矩阵切分到多块 GPU 上，每块 GPU 只计算矩阵乘法的一个``切片''：

\begin{figure}[H]
\centering
\includegraphics[width=0.95\textwidth]{slides-images/slide-040.jpg}
\caption{张量并行：将权重矩阵 $W$ 按列切分，每块 GPU 计算 $XW_i = Y_i$\protect\footnotemark}
\end{figure}
\footnotetext{Slides 第40页。}

\begin{knowledgebox}{两层 TP 的巧妙技巧}
对于 Transformer 的两层 MLP（$Y = \sigma(XW_1)W_2$），可以将 $W_1$ 按列切分、$W_2$ 按行切分。这样两层之间不需要通信，只在两层 MLP 结束后做一次 All-Reduce。这正好适配 Transformer 中 FFN 的两层 MLP 结构。
\end{knowledgebox}

\begin{figure}[H]
\centering
\includegraphics[width=0.95\textwidth]{slides-images/slide-041.jpg}
\caption{两层张量并行的数学原理：列切分第一层 + 行切分第二层 = 只需一次通信\protect\footnotemark}
\end{figure}
\footnotetext{Slides 第41页。}

\subsection{N 维并行：全部用上}

实际的大规模训练系统同时使用所有并行策略：

\begin{figure}[H]
\centering
\includegraphics[width=0.95\textwidth]{slides-images/slide-043.jpg}
\caption{Llama 3 405B 的 4D 并行配置：8-way TP $\times$ 16-way CP $\times$ 16-way PP $\times$ 8-way DP = 16,384 GPU\protect\footnotemark}
\end{figure}
\footnotetext{Slides 第43页。}

\begin{importantbox}{Llama 3 的 4D 并行}
在 16,384 块 GPU 的最大训练阶段：
\begin{itemize}
    \item \textbf{8-way 张量并行}：在单服务器的 8 块 GPU 上（最高带宽）
    \item \textbf{16-way 上下文并行}：处理 131K 长序列
    \item \textbf{16-way 流水线并行}：将 126 层分配给 16 个阶段
    \item \textbf{8-way 数据并行}：跨 Pod 复制（最低带宽）
\end{itemize}
关键原则：\textbf{高通信需求的并行策略放在高带宽连接上}。TP 通信最密集，放在服务器内部（900 GB/s）；DP 通信最稀疏，放在跨 Pod 连接上（50 GB/s）。
\end{importantbox}

\begin{figure}[H]
\centering
\includegraphics[width=0.95\textwidth]{slides-images/slide-044.jpg}
\caption{4D 并行的网络拓扑映射：不同并行策略对应不同的通信需求和网络层级\protect\footnotemark}
\end{figure}
\footnotetext{Slides 第44页。}

\subsection{本章小结}

上下文并行、流水线并行和张量并行分别从序列、层、维度三个方向切分计算。在实际的大规模训练中，四种并行策略同时使用，形成 4D 并行。不同策略的通信需求不同，应与集群的物理网络拓扑对齐。

%% ========== 容错 ==========

